\documentclass{article}
\usepackage{amsmath,hyperref}
\begin{document}
\title{Can a benchmark restore a resource while its fault is still active?}
\maketitle

\paragraph{Audit.}
For NIST ARIAC 2024.3.0, expiring one robot malfunction writes both
robot health flags true. Expiring one sensor blackout writes all six
sensor-type flags true. The writes do not check which other challenges
are still active. In the tested source prefix, a ceiling-robot stop can
restore the floor robot's healthy flag even though a separate floor-robot
stop is still live; the corresponding camera expiry also restores a
separately blacked-out logical camera flag. The switcher treats the
resulting floor health flag as a request to activate its controllers.
A replay of the pinned switcher also shows that a rejected switch is not
retried on repeated identical healthy reports. Both observations concern
benchmark software. Neither establishes an arm's physical state or
describes HeRo's implementation.

\paragraph{Source and cut.}
The \href{https://pages.nist.gov/ARIAC_docs/en/2024.1.0/competition/challenges.html}{published
rules} permit temporal robot malfunctions and sensor blackouts. They say
what the health flags represent and describe restoration when the
challenge resolves. The exact versions, paths and Git blob hashes are
in \texttt{SOURCE.tsv}. The task manager starts due challenges before
processing expiry. Each start updates the specified scope and retains an
occurrence object. Each expiry marks the occurrence completed, writes the
class-wide health reset, then erases an item from the same vector being
used by its range loop. An iterator after that erase cannot be relied on.
The harness therefore stops each selected native expiry function
immediately before its first erase and inspects the scalars and
occurrences at that cut. It passes those scalars to the unchanged health
message constructors. It does not run the original completion loop past
that cut, and makes no claim that a live Gazebo trial publishes the
captured state.

The first pair of temporal challenges disables the ceiling robot and the
ordinary camera on \([0,2)\). The second pair, when present, disables the
floor robot and the logical camera for five ticks, starting at tick 1
or tick 2. There is also a control with no second pair. At tick 2, the
first completion writes both robot and both selected camera flags true
in all three cases. The second pair's occurrence records still show it
started and not completed. The same-cut case tests the native start
before expiry ordering. The fixture pairs the robot and sensor
challenges on the same clock; a common physical cause is not asserted.

\paragraph{Downstream state.}
The pinned Python switcher is executed with controlled service and
subscription stubs. Feeding it those constructed robot messages requests
floor-group activation despite a live floor malfunction. With a
source-side union of live disabling scopes, the same switcher requests
no such activation in those histories and permits the healthy ceiling
group to resume. The no-second-pair control still permits floor service.
The logical camera plugin independently checks the logical camera flag
on its data path, so the flag is consumed as an enable condition there,
not proof that a fresh image has arrived. Its data path is source-audited
but not executed.

Even in the corrected health stream, a first floor activation request
can return \texttt{ok=false} without changing the mock plant. The native
switcher writes its local floor state while building the request; the
error branch logs but does not roll that state back. Three identical
later reports give no second request. A second stub check permits only
one controller to change even on a positive reply, illustrating the
limit of a group bit; it does not claim a positive real service reply
has that effect in this version. Recovery admission with a true health
flag and expected controllers should be followed by actual controller
membership feedback from an authoritative service before treating the
group as available. Current membership is not proof that a prior
physical effect has cleared. Waiting for an execution barrier or
retiring any live prior authority still belongs to the receiving
resource; the modeled stubs supply neither barrier nor clearance.

\paragraph{Composition law.}
For exact active occurrence set \(A\) and a typed disabling scope
\(S(e)\) for each occurrence, resource \(r\) is disabled just when
\[
  \exists e\in A:\ r\in S(e).
\]
An expiry of \(e\) removes only \(e\), not every cause with that fault
type. Unions across suppliers with the same occurrence identity
combine by disjunction on each resource; duplicate views of one
occurrence must not become new causes. This presupposes a common cut
and consistent definitions of shared occurrences. Independent sensor
views can disagree or be incomplete, in which case their apparent
union does not certify a release. At the authoritative event source
with exact transitions, induction over cuts proves the incremental
check equals the ordinary interval-union specification. The 57,624
pointwise comparisons cover every pair of nonempty masks on three
resources and every positive integer interval with start \(0\) to \(3\)
and end at most \(5\), including coincident starts and finishes.
Lean checks that another blocker survives erasure of one occurrence,
that union implements disjunction, and that resources outside all live
scopes need not stop.

\paragraph{Implication for evaluation.}
A test that scores recovery while the fault injector has ended a
different fault cannot certify concurrent tolerance merely from the
reported healthy flag. For the pinned benchmark, confirm the exact
injected active set at each cut or use an independently checked
occurrence-scoped injector, then distinguish controller state from
physical work. This does not say which releases the evaluated HeRo
scenarios used or whether they exercised either coincidence. A factored
direct checker given the same occurrence ledger agrees; the check is
a calibration for the integrated stack and admits activation
for an unaffected group in the declared fixture. It does not
demonstrate an advance over a correctly composed prior architecture or
discharge I1.
\end{document}
